\paragraph{The cost of family preprocessing}
This experiment tests whether family evidence accelerates the unchanged
catalogue. A useful policy outside it can coexist with false answers inside it.
We repeat the 240 queries on two 16-row jacket-type Fusion
views, $(\Lambda,\rho)=(1.5,0)$ and $(2,0)$, selected retrospectively with known
capacities. Hot-start (Hot), optimized threshold execution (Opt), and Opt with
partition (Part) or exact-LP (LP) preprocessing share objective caching, checked
points and a rational receiver. Family certificates add only feasible points;
a shared constant-action shortcut skips unnecessary family solves.
Unresolved represented queries permit one producer attempt per objective--margin
pair and one complete exact fallback per objective; ideal answers reuse this
evidence without extra solves.

\begin{table}[t]
\centering
\caption{Original Coat queries. Time: median seconds over three rounds, including
startup, family work, external verification and cleanup. Ratio/Opt: paired
cost range (above one is slower). N/E/F: native/fixed-exact/family calls,
identical per stream across quality and timing. Both views use $\rho=0$.}
\label{tab:coat-workflow}
\begingroup\small
\renewcommand{\arraystretch}{1.10}\setlength{\tabcolsep}{4pt}
\begin{tabular}{@{}llrrr@{}}
\toprule
\rowcolor{QHeader}\qth{View} & \qth{Plan} & \qth{Time (s)} & \qth{Ratio/Opt} & \qth{N/E/F} \\
\midrule
$\Lambda=1.5$ & Hot & 6.800 & 0.974--1.009 & 2/2/0 \\
$\Lambda=1.5$ & Opt & 6.934 & 1.000 & 2/2/0 \\
$\Lambda=1.5$ & Opt+Part & 7.312 & 1.051--1.077 & 2/2/1 \\
$\Lambda=1.5$ & Opt+LP & 7.276 & 1.042--1.079 & 2/2/1 \\
\midrule
$\Lambda=2$ & Hot & 6.329 & 0.985--0.998 & 1/1/0 \\
$\Lambda=2$ & Opt & 6.390 & 1.000 & 1/1/0 \\
$\Lambda=2$ & Opt+Part & 6.431 & 0.999--1.024 & 1/1/0 \\
$\Lambda=2$ & Opt+LP & 6.359 & 0.990--1.002 & 1/1/0 \\
\bottomrule
\end{tabular}\endgroup

\end{table}

All 32 streams complete, each with 240 false answers under both semantics.
On the nonzero-capacity view, Part adds 5.1--7.7\% to Opt's full cost without
reducing its two native and two exact calls (Table~\ref{tab:coat-workflow});
LP likewise saves no calls. The zero-capacity control skips family computation.
Three rotations describe machine variation; execution positions are not fully
balanced. Different optimal points compare complete preprocessing choices.
Family answers enable redesign beyond the catalogue but do not amortize
preprocessing here.
